% TOP_PROBLEM: {"id": 3116, "title": "Star chromatic index of complete graphs", "category_id": 3, "category": {"id": 3, "name": "graph_theory", "display_name": "Graph Theory", "description": "Problems involving graphs, networks, and their properties.", "slug": "graph-theory", "order_index": 3, "created_at": "2026-07-31T15:26:25.671Z"}, "status": "open", "problem_number": "OPG-37275", "set_id": 12}
% FIRST_POSTED: 2026-10-01
% ENTRY_KIND: statement_only
% UnsolvedMath ID 3116; source catalogue code OPG-37275
% !TeX program = lualatex
% Definition and sources only; this file is not an LLM solution attempt.
\documentclass[11pt,a4paper]{article}
\usepackage[margin=25mm]{geometry}
\usepackage{fontspec}
\setmainfont{lmroman10-regular.otf}[BoldFont=lmroman10-bold.otf,
 ItalicFont=lmroman10-italic.otf,BoldItalicFont=lmroman10-bolditalic.otf]
\usepackage{amsmath,amssymb,mathtools}
\usepackage{unicode-math}
\setmathfont{latinmodern-math.otf}
\AtBeginDocument{\let\setminus\smallsetminus}
\usepackage{xurl,hyperref}
\hypersetup{colorlinks=true,urlcolor=blue}
\setcounter{secnumdepth}{0}
\setlength{\parindent}{0pt}
\setlength{\parskip}{5pt}
\setlength{\emergencystretch}{3em}
\begin{document}
\section*{Star chromatic index of complete graphs}\label{3116}
{\small UnsolvedMath ID 3116 · OPG-37275 · Graph Theory · OpenGarden}
\subsection{Definitions and mathematical statement}
For $n\ge2$, let $K_n$ be the complete simple undirected graph on
$\{1,\ldots,n\}$. A proper edge colouring assigns different colours to
any two edges sharing a vertex. It is a star edge colouring if, in
addition, each path with four edges and each cycle with four edges
receives at least three colours. Paths are simple, with five distinct
vertices, and need not be induced; this distinction matters in a
complete graph.

Let $\chi'_s(K_n)$ denote the minimum size of a colour palette admitting
such a colouring. The question is whether there exists an absolute
constant $C>0$ such that
\[
 \chi'_s(K_n)\le Cn\qquad\text{for every }n\ge2.
\]
This is the precise meaning of the requested $O(n)$ upper bound.
Allowing the inequality only for all sufficiently large $n$ is equivalent,
since the finitely many smaller values can be absorbed by increasing $C$.

The palette and colouring may change with $n$, but $C$ must not grow
with $n$. Every pair of colours induces a union of paths of at most
three edges, which gives another way to express the forbidden-pattern
condition. Since all $n-1$ edges at a vertex require different colours,
a linear upper bound would also give the order of growth
$\chi'_s(K_n)=\Theta(n)$.
The problem does not ask that the same constant work for an arbitrary
graph in terms of its maximum degree; it is the family of complete
graphs that is quantified over.
\subsection{Short English statement}
Is a constant times $n$ colours enough to star edge-colour every complete graph $K_n$?
\subsection{Sources}
\begin{itemize}
\item[S1] ULAM.AI and the UnsolvedMath Contributors, supplied problems.json, record 3116 (OPG-37275), OpenGarden collection. Catalogue identity and source transcription; CC BY 4.0.
\url{https://huggingface.co/datasets/ulamai/UnsolvedMath}.
\item[S2] Open Problem Garden, Star chromatic index of complete graphs. Original problem and discussion; the mathematical formulation below is expanded from the supplied source record.
\url{http://www.openproblemgarden.org/op/star_chromatic_index_of_complete_graphs}.
\end{itemize}
\end{document}
